% Figure: effectiveness as a function of NTU for a counter-flow exchanger at
% several values of the heat-capacity-rate ratio C_r, plus the C_r=0 phase-change
% limit. The chart is the working design tool of the NTU method.
\begin{figure}[htbp]
\centering
\begin{tikzpicture}
\begin{axis}[
  width=12cm, height=8cm,
  xlabel={number of transfer units, $\mathrm{NTU}=UA/C_{\min}$},
  ylabel={effectiveness $\varepsilon$},
  xmin=0, xmax=5, ymin=0, ymax=1.0,
  axis lines=left,
  legend pos=south east,
  legend cell align=left,
  tick label style={font=\scriptsize},
  label style={font=\small},
  legend style={font=\scriptsize},
  samples=80,
  domain=0.001:5,
]
% C_r = 0 (phase change): eps = 1 - exp(-NTU)
\addplot[engblue, very thick] {1 - exp(-x)};
\addlegendentry{$C_r=0$ (condenser/evaporator)}
% C_r = 0.5 counter-flow
\addplot[engamber, very thick]
  {(1 - exp(-x*(1-0.5)))/(1 - 0.5*exp(-x*(1-0.5)))};
\addlegendentry{$C_r=0.5$ counter-flow}
% C_r = 1 counter-flow: eps = NTU/(1+NTU)
\addplot[engred, very thick] {x/(1+x)};
\addlegendentry{$C_r=1$ counter-flow}
\end{axis}
\end{tikzpicture}
\caption[Effectiveness versus NTU for counter-flow]{Effectiveness as a function
of the number of transfer units for a counter-flow exchanger at three values of
the heat-capacity-rate ratio $C_r=C_{\min}/C_{\max}$. The top curve is the
$C_r=0$ limit, which holds for a condenser or evaporator where one stream changes
phase at constant temperature; there the effectiveness reaches $1-e^{-\mathrm{NTU}}$
independent of geometry. As $C_r$ rises toward unity the effectiveness for a
given area falls, because the second stream warms (or cools) appreciably and
erodes the driving difference. The diminishing slope past $\mathrm{NTU}\approx 3$
is the economic ceiling of the method: doubling the area near the knee buys very
little extra effectiveness, which is why approach temperatures below a few kelvin
are expensive.}
\label{fig:vol04ch06:effectiveness-ntu}
\end{figure}
